\chapter{Operations}\label{ch:fm:operations}

On 28 May 2024 most US securities trades moved from settlement two business days after the trade to one. A firm that had had a day to find a
missing settlement instruction now had the evening of the trade date: a new rule required allocations, confirmations and affirmations to be
completed by the end of the trade date, and the industry's affirmation cut-off was 9 p.m. Eastern time. The industry report that followed found
that nearly 95\% of transactions met that cut-off, against 73\% in January 2024, and that fail rates stayed where they had been under the two-day
cycle. The work had not disappeared. It had moved into fewer hours, and every operations team that did not want more fails had to be sized for the
shorter window.

\section{Front, middle and back office}

A trade is a promise to exchange securities for cash on a later date. Between the promise and the exchange lie the steps of Book 1's settlement
chapter and of Book 15's post-trade platform: trade capture, enrichment, confirmation, clearing, settlement instruction, settlement, and the
reconciliation of every record with every counterparty's. A trading firm divides the work into three groups.

\begin{definition}[Front, middle and back office]\label{def:fm:operations:offices}
The \emph{front office}\index{front office} makes and executes the trading decisions: traders, researchers and sales. The \emph{middle
office}\index{middle office} controls and reports on what the front office did, independently of it: trade validation, P\&L and position
reporting, risk control and product control. The \emph{back office}\index{back office} processes the trades to completion: confirmation,
clearing, settlement, cash and collateral movements, corporate actions and reconciliation.
\end{definition}

The division is also a control. Segregation of duties (Book 6) requires that the person who makes a trade does not book, confirm or settle it,
and that the person who reports its P\&L does not trade it: the \omterm{def:fm:operations:offices}{middle office} exists so that the \omterm{def:fm:operations:offices}{front office}'s numbers are checked by someone who
is not paid on them. Operational risk (Book 6) lives in the hand-offs between the three.

\begin{omfigure}
\begin{tikzpicture}[font=\footnotesize,>=stealth,box/.style={draw,rounded corners,align=center,minimum height=0.75cm,text width=2.2cm}]
\node[box,fill=omDef!15] (fo) at (0,0) {front office\\{\scriptsize execute}};
\node[box,fill=omProp!15] (mo) at (3.1,0) {middle office\\{\scriptsize validate, report}};
\node[box,fill=omThm!15] (bo) at (6.2,0) {back office\\{\scriptsize confirm, settle}};
\node[box,fill=gray!15] (st) at (9.3,0) {settlement\\{\scriptsize securities for cash}};
\draw[->] (fo) -- (mo); \draw[->] (mo) -- (bo); \draw[->] (bo) -- (st);
\draw[thick] (0,-1.3) -- (9.3,-1.3);
\foreach \x/\l in {0/T,4.65/{T+1},9.3/{T+2}} {\draw (\x,-1.2) -- (\x,-1.4) node[below] {\l};}
\draw[<->,omThm,thick] (0,-2.2) -- node[below,font=\scriptsize] {two-day cycle: fix breaks during T+1} (9.3,-2.2);
\draw[<->,omDef,thick] (0,-2.9) -- node[below,font=\scriptsize] {one-day cycle: affirm by the evening of T} (4.65,-2.9);
\end{tikzpicture}

\omcaption{The three offices along a trade's life, and the window in which a break must be fixed before settlement under a two-day and a one-day
cycle. Schematic.}\label{fig:fm:operations:offices}
\end{omfigure}

\begin{dated}{2026-09}{Settlement cycles}\label{dat:fm:operations:cycle}
\textbf{United States}: SEC Release 34-96930 (February 2023) amended Rule 15c6-1 so that a broker-dealer may not, for most securities, agree to
settle later than the first business day after the trade unless the parties expressly agree otherwise at the time of the trade; compliance began
on 28 May 2024. New Rule 15c6-2 requires allocations, confirmations and affirmations to be completed as soon as technologically practicable and no
later than the end of the trade date. SIFMA, ICI and DTCC reported in September 2024 that nearly 95\% of transactions met the 9 p.m. ET affirmation
cut-off (73\% in January 2024) and that the July 2024 fail rates, 2.12\% in the continuous net settlement system and 3.31\% outside it, were
consistent with those under T+2. \textbf{European Union}: Regulation (EU) 2025/2075 moves the settlement cycle of the Central Securities
Depositories Regulation from two business days to one from 11 October 2027.
\end{dated}

\section{Reconciliation and breaks}

A firm's records of a trade, a position or a cash balance must agree with its counterparty's, its clearing house's, its \omterm{def:fm:the-asset-manager-and-the-fund:admin}{custodian}'s and its own
other systems. Reconciliation compares them; a difference is a reconciliation break (Book 15), and its age since it was found is the break ageing
that the operations team reports every day. Breaks have causes: a wrong price or quantity at capture, a missing or stale settlement instruction, an
allocation not sent, a corporate action booked differently. Most are resolved in minutes; a few take days; a break still open when the securities
are due to move becomes a settlement fail.

The team that resolves them is a queue (Book 4): breaks arrive, wait for someone free and are worked. Let breaks arrive at rate $\lambda$ an hour,
let each take an exponential time with mean $1/\mu$ to resolve, and let $c$ people work them first come, first served. The queue is stable if
$\lambda<c\mu$. A break is late if its time in the system, wait plus work, exceeds the time $D$ left before settlement.

\begin{proposition}[Late breaks in a team of $c$]\label{prop:fm:operations:late}
In the stationary M/M/c queue with $a=c\mu-\lambda>0$, let $C$ be the probability that an arriving break waits (Erlang's C formula). The
probability that a break is not resolved within $D$ is
\[
P(T>D)=(1-C)\,e^{-\mu D}+C\,\frac{a\,e^{-\mu D}-\mu\,e^{-aD}}{a-\mu}\qquad(a\neq\mu),
\]
and $P(T>D)\ge e^{-\mu D}$ for every $c$: if $e^{-\mu D}>1-s$, no team is large enough to resolve a share $s$ of breaks in time.
\end{proposition}

\begin{proof}[Partial proof]
In the FCFS M/M/c queue an arriving break waits with probability $C$, and given that it waits, its wait is exponential with rate $a$ (a standard
result, \admitted). Its work $S$ is exponential with rate $\mu$ and independent of the wait. With no wait, $P(S>D)=e^{-\mu D}$; with a wait $W$,
$P(W+S>D)$ is the tail of the sum of two independent exponentials, which is the fraction shown. Both terms are at least $e^{-\mu D}$, since
$W\ge0$.
\end{proof}

\omcode{code/firm/opsmetrics/firm_opsmetrics.py}{39}{68}{The waiting probability from firm.queues' stationary law, the late share of
\cref{prop:fm:operations:late}, and the smallest team for a service level.}\label{lst:fm:operations:late}

The floor is the point of the proposition. Staff shorten the wait, not the work: once no break waits, the late share is the share of breaks whose
work alone takes longer than the window. Under a two-day cycle the window is long and the floor invisible; under a one-day cycle it can be the
binding constraint.

\section{Tutorial: one day less}\label{tut:fm:operations:tut}

\begin{tutorial}
\textbf{Goal.} Size a breaks team for a two-day and a one-day cycle, and price the fails it lets through. \textbf{End state:}
\cref{fig:fm:operations:late}, \cref{fig:fm:operations:cost} and the staff needed for a 95\% service level.
\begin{tutsteps}
\item \textbf{The flow.} 120 breaks a day arrive over an eight-hour day ($\lambda=15$ an hour) and take 40 minutes of work on average ($\mu=1.5$);
the offered load is $\lambda/\mu=10$ people (illustrative).
\item \textbf{The windows.} A break must be resolved within 10 working hours under T+2 and within 2.5 under T+1, before the evening cut-off.
\item \textbf{The fails.} Half the late breaks become fails, on trades of \$5 million, failing for two days at a penalty of 1 basis point a day
(the EU rate for liquid shares, an input) plus \$1\,500 of internal cost: \$2\,500 a fail. A person costs \$120\,000 a year.
\item \textbf{The sizing.} \texttt{fm\_ops.staff95(cycle)} applies \texttt{firm.opsmetrics.staff\_for}; \texttt{fm\_ops.table(cycle)} prices
teams by staff and fail cost.
\end{tutsteps}
\end{tutorial}

\begin{omfigure}
\begin{tikzpicture}
\begin{axis}[width=10.5cm,height=5.2cm,xlabel={\small people in the team},ylabel={\small breaks late (\%)},xmin=11,xmax=20,ymin=0,ymax=10,
  tick label style={font=\footnotesize},grid=major,grid style={gray!20},table/col sep=comma,
  legend cell align=left,legend style={font=\scriptsize,at={(0.5,-0.3)},anchor=north,/tikz/every even column/.append style={column sep=8pt}},legend columns=3]
\addplot[omThm,thick,mark=*,mark size=1.2pt] table[x=team,y=d10]{figdata/desk/13-operations/late.csv};
\addplot[omProp,thick,dashed,mark=square*,mark size=1.2pt] table[x=team,y=d4]{figdata/desk/13-operations/late.csv};
\addplot[omDef,thick,mark=triangle*,mark size=1.5pt] table[x=team,y=d25]{figdata/desk/13-operations/late.csv};
\draw[gray,dotted] (axis cs:11,5) -- (axis cs:20,5);
\draw[omDef,densely dotted] (axis cs:11,2.35) -- (axis cs:20,2.35);
\legend{10-hour window (T+2),4-hour window,2.5-hour window (T+1)}
\end{axis}
\end{tikzpicture}

\omcaption{The share of breaks not resolved within the settlement window, by team size, for 120 breaks a day of 40 minutes' work. With 2.5 hours no
team gets below the floor $e^{-\mu D}=2.35\%$ (lower dotted line) set by the work itself; the upper dotted line is the 5\% of a 95\% service level.
Data: \texttt{firm.opsmetrics.p\_late}.}\label{fig:fm:operations:late}
\end{omfigure}

Under T+2 eleven people, the smallest stable team, resolve all but a few breaks in a million within the window: the late share is 0.0003\% and
the year's fail cost is nothing. Under T+1 eleven people let 8.37\% of breaks run late, 1\,255 fails a year costing \$3.14 million; twelve are needed
for the 95\% service level (3.38\% late, \$1.27 million) and thirteen minimise staff and fail cost together (2.69\% late, \$2.57 million in all,
against \$1.32 million under T+2). Beyond thirteen, people buy almost nothing: the late share approaches the 2.35\% floor
(\cref{fig:fm:operations:cost}). Waiting is the smaller part of the problem: with thirteen people 28.5\% of breaks wait at all.

\begin{omfigure}
\begin{tikzpicture}
\begin{axis}[width=10.5cm,height=5.2cm,xlabel={\small people in the team},ylabel={\small \$ million a year},xmin=11,xmax=18,ymin=0,ymax=5,
  tick label style={font=\footnotesize},grid=major,grid style={gray!20},table/col sep=comma,
  legend cell align=left,legend style={font=\scriptsize,at={(0.5,-0.3)},anchor=north,/tikz/every even column/.append style={column sep=8pt}},legend columns=3]
\addplot[omThm,thick,dashed] table[x=team,y=staff]{figdata/desk/13-operations/cost.csv};
\addplot[omProp,thick,dotted] table[x=team,y=fail]{figdata/desk/13-operations/cost.csv};
\addplot[omDef,very thick,mark=*,mark size=1.2pt] table[x=team,y=total]{figdata/desk/13-operations/cost.csv};
\legend{staff cost,fail cost,total}
\end{axis}
\end{tikzpicture}

\omcaption{Staff cost, fail cost and their sum under a one-day cycle. The total is lowest at thirteen people, \$2.57 million a year, where the fail
cost is still \$1.01 million: the floor of \cref{prop:fm:operations:late} cannot be staffed away. Data: \texttt{fm\_ops.table}.}\label{fig:fm:operations:cost}
\end{omfigure}

\begin{table}[ht]
\centering\small
\begin{tabular}{lrrrrr}
\toprule
scenario & team for 95\% & cheapest team & staff cost & fail cost & total \\
\midrule
T+2 & 11 & 11 & 1.32 & 0.00 & 1.32 \\
T+1 & 12 & 13 & 1.56 & 1.01 & 2.57 \\
T+1, half the breaks & 7 & 7 & 0.84 & 0.58 & 1.42 \\
T+1, 30-minute work & 9 & 10 & 1.20 & 0.30 & 1.50 \\
\bottomrule
\end{tabular}
\caption{Staffing a breaks team under each cycle, and two ways of changing the work itself ($\$$ million a year). Data:
\texttt{fm\_ops.scenarios}.}\label{tab:fm:operations:scenarios}
\end{table}

The remedy is to change the work, not the team (\cref{tab:fm:operations:scenarios}). Halving the breaks, by automating the steps that cause
them (standing settlement instructions, electronic allocation and matching on the trade date), needs seven people and costs \$1.42 million; making
the work faster, at 30 minutes a break, needs ten and costs \$1.50 million. Either almost undoes the cost of the shorter cycle, and staffing alone
cannot.

\begin{remark}[What the queue leaves out]\label{rem:fm:operations:batch}
Breaks do not arrive evenly: many appear together, when the end-of-day reconciliations run, which is the worst time under a one-day cycle. A queue
with batch arrivals at the close gives larger late shares than the chapter's even flow, and makes the case for moving reconciliation to real time
(Book 15) rather than for more people in the evening.
\end{remark}

\begin{omfigure}
\begin{tikzpicture}
\begin{axis}[width=10.5cm,height=5cm,ybar,bar width=8pt,ylabel={\small share of breaks (\%)},ymin=0,ymax=60,xtick={0,1,2,3,4},
  xticklabels={{$<$30 min},{30--60 min},{1--2 h},{2--4 h},{$>$4 h}},xmin=-0.6,xmax=4.6,
  tick label style={font=\footnotesize},grid=major,grid style={gray!20},table/col sep=comma,
  legend cell align=left,legend style={font=\scriptsize,at={(0.5,-0.25)},anchor=north,/tikz/every even column/.append style={column sep=8pt}},legend columns=3]
\addplot[fill=omDef!55,draw=omDef] table[x=pos,y=c11]{figdata/desk/13-operations/ageing.csv};
\addplot[fill=omProp!55,draw=omProp] table[x=pos,y=c12]{figdata/desk/13-operations/ageing.csv};
\addplot[fill=omThm!50,draw=omThm] table[x=pos,y=c14]{figdata/desk/13-operations/ageing.csv};
\legend{11 people,12 people,14 people}
\end{axis}
\end{tikzpicture}

\omcaption{The ageing profile of resolved breaks: how long each took from arrival to resolution, for three team sizes. More people move breaks
into the first half hour; the share older than four hours falls from 1.3\% to 0.3\%, but not to zero. Data: \texttt{fm\_ops.ageing}.}\label{fig:fm:operations:ageing}
\end{omfigure}

\section{Corporate actions and reference events}

A dividend, a split, a merger, a rights issue or a coupon changes what a position is worth or what it is. Book 1 describes the events; for
operations they are breaks waiting to happen, since every system that holds the position must apply the event the same way on the same date, and
the firm must decide on elective events (take the cash or the shares, tender or not) before a deadline set by the issuer's agent. The controls are
those of reconciliation, applied before the event rather than after it: a golden source of reference data, an event calendar checked against the
positions that will be affected, and a named owner for every elective decision.

\section{Collateral management}

Every margined relationship moves collateral daily: variation margin with clearing houses (Book 1) and, under a credit support annex (Book 2), with
bilateral counterparties above a threshold and in amounts no smaller than the minimum transfer amount (Book 6).

\begin{definition}[Collateral management, collateral dispute]\label{def:fm:operations:collateral}
\emph{Collateral management}\index{collateral management} is the daily process of computing, calling, paying, receiving, valuing and substituting
the collateral due under a firm's margin agreements, and of reconciling its records with its counterparties'. A \emph{collateral
dispute}\index{collateral dispute} arises when the two parties' calculations of the amount due differ by more than an agreed tolerance; the
undisputed amount is usually paid while the difference is investigated.
\end{definition}

\begin{example}[A disputed call]\label{ex:fm:operations:call}
A CSA has a threshold of \$10 million and a minimum transfer amount of \$0.5 million; the firm holds \$14.6 million. The firm values the exposure
at \$25.3 million and calls $25.3-10-14.6=\$0.7$ million. The counterparty values it at \$24.4 million: its figure is a return of \$0.2 million,
below the minimum transfer amount, so it sees no call. The valuations differ by 3.6\%, above a 2\% tolerance: the call is disputed, and the cause
is almost always a difference in trade population or in the marks of a few trades, which is where the reconciliation starts.
\end{example}

\omcode{code/firm/opsmetrics/firm_opsmetrics.py}{81}{106}{A collateral call under a CSA's threshold and minimum transfer amount, the dispute
test, and a key risk indicator's traffic light.}\label{lst:fm:operations:csa}

\section{Settlement discipline}

\begin{definition}[Settlement discipline regime]\label{def:fm:operations:discipline}
A \emph{settlement discipline regime}\index{settlement discipline regime} is the set of rules by which a market deters and handles settlement fails:
the reporting of fails, cash penalties charged to the failing party, and buy-ins that replace securities not delivered.
\end{definition}

\begin{dated}{2026-09}{The EU settlement discipline regime}\label{dat:fm:operations:csdr}
The Central Securities Depositories Regulation's settlement discipline rules (Delegated Regulation (EU) 2018/1229) have applied since 1 February
2022. Cash penalties are charged per day of fail on the value of the failed instruction, at rates set in Delegated Regulation (EU) 2017/389: 1.0
basis point for liquid shares, 0.5 for illiquid shares, 0.25 for SME growth market instruments, 0.10 for sovereign debt, 0.20 for other debt,
0.15 for SME growth market debt and 0.5 for other instruments; a lack of cash is charged at the issuing central bank's overnight credit rate,
floored at zero. Regulation (EU) 2023/2845 made mandatory buy-ins a measure of last resort: they apply only if the Commission introduces them by
an implementing act, after a cost-benefit analysis by ESMA, where other measures have not reduced fails.
\end{dated}

A penalty of 1 basis point a day on a \$5 million trade is \$500 a day: small against a trading desk's P\&L, but charged on every fail every day,
and paid to the counterparty that did not fail. The larger costs of a fail are those the tutorial prices as internal: funding the unsettled
position, claims, the time spent, and the damage to the relationship. A \omterm{def:fm:operations:kri}{key risk indicator} makes them visible before they happen.

\begin{definition}[Key risk indicator]\label{def:fm:operations:kri}
A \emph{key risk indicator}\index{key risk indicator} is a metric that is monitored against thresholds because it moves before losses do: for
operations, the number of breaks older than a day, the fail rate, the share of trades affirmed on the trade date, the number of disputed calls.
\end{definition}

\begin{method}[A daily operations report]\label{met:fm:operations:report}
\begin{enumerate}
\item Choose five to ten \omterm{def:fm:operations:kri}{key risk indicators}, each with an amber and a red threshold and an owner.
\item Compute them from the systems of record, not from the teams' own counts.
\item Report the traffic lights daily to the \omterm{def:fm:operations:offices}{middle office} and the head of operations, and weekly with trends to the \omterm{def:fm:the-risk-management-function:cro}{risk committee}.
\item For every red, record the cause and the fix; for every recurring amber, change the process, not the threshold.
\end{enumerate}
\end{method}

On the chapter's report (\texttt{fm\_ops.daily\_report}), fourteen breaks older than a day, a fail rate of 2.4\% and 93\% of trades affirmed on the
trade date are all amber against thresholds of 10, 2\% and 95\%; one disputed call is green.

\section{Build: operations metrics}\label{bld:fm:operations:opsmetrics}

\begin{build}
\bfield{Purpose} The operations team's numbers: the break queue and its ageing, the staff for a service level, fail costs, collateral calls and
disputes, and \omterm{def:fm:operations:kri}{key risk indicators}.
\bfield{Interface} \texttt{firm.opsmetrics}: \texttt{mmc\_stationary}, \texttt{p\_wait}, \texttt{erlang\_c}, \texttt{p\_late}, \texttt{staff\_for},
\texttt{ageing}, \texttt{fail\_cost}; \texttt{Csa}, \texttt{call}, \texttt{disputed}; \texttt{kri}, \texttt{report}.
\bfield{Rules} The queue must be stable; a service level below the work-time floor is refused; penalty rates are inputs, never constants; a \omterm{def:fm:operations:kri}{key risk
indicator} whose red lies below its amber is read the other way.
\bfield{Acceptance tests} \texttt{code/firm/opsmetrics/tests/}: the chain's waiting probability against Erlang's formula; the floor; staffing at the
boundary; calls, disputes and traffic lights on hand numbers.
\bfield{Stretch} Batch arrivals at the close; break types with their own work times; a collateral workflow with substitution and ageing of
disputes.
\end{build}

\begin{omsources}
\item SEC Release 34-96930 (2023); 17 CFR 240.15c6-1 and 240.15c6-2.
\item SIFMA, ICI and DTCC, T+1 After Action Report, press release of 12 September 2024.
\item Delegated Regulations (EU) 2017/389 and 2021/70; Regulations (EU) 2023/2845 and 2025/2075.
\end{omsources}

\section{Exercises}

\begin{exercise}[$\star$]\label{exo:fm:operations:1}
Which office does each of these belong to: booking a correction to a trade's price, approving the day's P\&L, choosing the cash option of a
dividend, executing an order?
\end{exercise}

\begin{exercise}[$\star$]\label{exo:fm:operations:2}
What does one fail cost under the tutorial's assumptions, and how much of it is the penalty?
\end{exercise}

\begin{exercise}[$\star$]\label{exo:fm:operations:3}
With breaks taking 40 minutes on average, what is the floor on the late share for a 2.5-hour window and for a 10-hour window?
\end{exercise}

\begin{exercise}[$\star\star$]\label{exo:fm:operations:4}
Compute the probability that a break waits with 11, 12 and 14 people ($\lambda=15$, $\mu=1.5$).
\end{exercise}

\begin{exercise}[$\star\star$]\label{exo:fm:operations:5}
In \cref{ex:fm:operations:call}, what does each party think is due, and why is the call disputed?
\end{exercise}

\begin{exercise}[$\star\star$]\label{exo:fm:operations:6}
Why does adding people beyond thirteen barely reduce fails under T+1, and what does?
\end{exercise}

\begin{exercise}[$\star\star\star$]\label{exo:fm:operations:7}
\emph{Coding.} Run \texttt{fm\_ops.scenarios()}. Which change to the work saves more, halving the breaks or cutting the work to 30 minutes, and
why does the faster work need more people?
\end{exercise}

\begin{exercise}[$\star\star\star$]\label{exo:fm:operations:8}
\emph{Find the flaw.} ``The move to T+1 halves our window, so we will double the operations team.''
\end{exercise}

\section{Problem: One Day Less}

\begin{problem}[{Weekend problem --- one day less}]\label{pb:fm:operations:1}
A head of operations must present a plan for a firm's move from a two-day to a one-day settlement cycle.

\medskip\noindent\textbf{Part I --- The organisation.}
\begin{enumerate}
\item Define the front, middle and \omterm{def:fm:operations:offices}{back office}.
\item Why is the \omterm{def:fm:operations:offices}{middle office} independent of the \omterm{def:fm:operations:offices}{front office}?
\item State the US rules on the settlement cycle and on same-day affirmation.
\item What did the industry report find after the US move?
\end{enumerate}

\medskip\noindent\textbf{Part II --- The queue.}
\begin{enumerate}[resume]
\item Model the breaks team as a queue and state its stability condition.
\item State \cref{prop:fm:operations:late} and prove the floor.
\item Give the floor for the tutorial's T+1 window.
\item Give the probability that a break waits with 11, 12 and 14 people.
\end{enumerate}

\medskip\noindent\textbf{Part III --- The sizing.}
\begin{enumerate}[resume]
\item Give the staff needed for a 95\% service level under each cycle.
\item Give the late share, fails and fail cost of eleven, twelve and thirteen people under T+1.
\item Which team minimises the total cost under T+1, and what is that cost against T+2?
\item Give the cost of the two changes to the work in \cref{tab:fm:operations:scenarios}.
\item How does the ageing profile change from eleven to fourteen people?
\item What does the model leave out, and in which direction does that bias it?
\end{enumerate}

\medskip\noindent\textbf{Part IV --- Discipline and collateral.}
\begin{enumerate}[resume]
\item Define a \omterm{def:fm:operations:discipline}{settlement discipline regime} and give the EU penalty rates for shares and sovereign debt.
\item What changed for buy-ins in 2023?
\item Define \omterm{def:fm:operations:collateral}{collateral management} and a \omterm{def:fm:operations:collateral}{collateral dispute}, and work through \cref{ex:fm:operations:call}.
\item Define a \omterm{def:fm:operations:kri}{key risk indicator} and give four for operations.
\item State the \emph{named result}: the staff needed to keep 95\% of breaks resolved before settlement under a two-day and a one-day cycle, and
the fail cost avoided.
\item In two sentences, write the plan.
\end{enumerate}
\end{problem}

\section{Interview questions}

\begin{interviewq}[$\star$ \iqroles{trader, developer}]\label{iq:fm:operations:1}
What is the difference between the \omterm{def:fm:operations:offices}{middle office} and the \omterm{def:fm:operations:offices}{back office}?
\end{interviewq}

\begin{interviewq}[$\star$ \iqroles{developer}]\label{iq:fm:operations:2}
What is a reconciliation break, and what are its usual causes?
\end{interviewq}

\begin{interviewq}[$\star\star$ \iqroles{researcher}]\label{iq:fm:operations:3}
Breaks arrive at 15 an hour and take 40 minutes each. How many people do you need at the least, and why is that not enough?
\end{interviewq}

\begin{interviewq}[$\star\star$ \iqroles{risk}]\label{iq:fm:operations:4}
A counterparty disputes your margin call. What do you do?
\end{interviewq}

\begin{interviewq}[$\star\star$ \iqroles{trader, risk}]\label{iq:fm:operations:5}
What did the US move to T+1 change for a trading firm's operations?
\end{interviewq}

\begin{interviewq}[$\star\star\star$ \iqroles{developer, researcher}]\label{iq:fm:operations:6}
Design a system that reduces settlement fails under a one-day cycle.
\end{interviewq}
